\documentclass[
    11pt,
    a4paper,
    egregdoesnotlikesansseriftitles,
    toc=chapterentrywithdots,
    twoside,openright,
    titlepage,
    parskip=half,
    headings=normal,  % reduces heading size
    listof=totoc,
    bibliography=totoc,
    index=totoc,
    captions=tableheading,  % caption below table
    chapterprefix,
    listof=flat,
    final
]{scrbook}


% details about your thesis
\newcommand{\titel}{My very fancy thesis title.}
\newcommand{\artderarbeit}{Bachelorarbeit}  % {Bachelorarbeit,Masterarbeit}
\newcommand{\autor}{Georg Simon Ohm}
\newcommand{\studiengang}{Informatik}  % {Informatik,Wirtschaftsinformatik,Medieninformatik}
\newcommand{\matrikelnr}{123\,4567}
\newcommand{\erstgutachter}{Prof.\,Dr.~Korbinian Riedhammer}
\newcommand{\zweitgutachter}{Prof.\,Dr.~Bartosz von\,Rymon\,Lipinski}
\newcommand{\betreuer}{M.Sc.\,~Martina Schmidt}
\newcommand{\unternehmen}{Musterfirma GmbH}
\newcommand{\logo}{figures/TH-Nuernberg-RGB.png}
\newcommand{\keywords}{hot, fuzz}
 

% custom head and foot
\usepackage[automark]{scrlayer-scrpage}
\pagestyle{scrheadings}
\ihead{\headmark}
\chead{}
\ohead{\pagemark}
\renewcommand*\chaptermarkformat{\chapappifchapterprefix{\ }% 
  \thechapter.\enskip}

\RedeclareSectionCommand[tocindent=0pt]{section}
\RedeclareSectionCommand[tocindent=0pt]{subsection}
%\RedeclareSectionCommand[tocnumwidth=70pt]{chapter}


%% Farben (Angabe in HTML-Notation mit großen Buchstaben)

\usepackage{xcolor} 	% für HTML-Notation

\definecolor{LinkColor}{HTML}{00007A}
\definecolor{ListingBackground}{HTML}{FAFAFA}
\definecolor{keywords}{rgb}{0.133,0.133,0.6}
\definecolor{comments}{HTML}{848484}
\definecolor{strings}{rgb}{0.133,0.545,0.133}
\colorlet{punct}{red!60!black}
\definecolor{delim}{RGB}{20,105,176}
\colorlet{numb}{magenta!60!black}

\usepackage{scrhack}

% other packages
\usepackage[utf8]{inputenc}
\usepackage[T1]{fontenc}
\usepackage{lmodern,relsize,textcomp,csquotes}
\usepackage{amsmath,amsfonts}
\usepackage[english, ngerman]{babel}  % flip for english thesis
\usepackage[final]{graphicx}
\usepackage{setspace,geometry,xcolor}
\usepackage{makeidx}
\usepackage{paralist,ifthen,todonotes}
\usepackage{url}
\usepackage[toc]{glossaries}
\usepackage{pdfpages}

% table setup
\usepackage{longtable}
\usepackage{array}
\usepackage{ragged2e}
\usepackage{lscape}
\usepackage{tabularx}
\newcolumntype{L}[1]{>{\raggedright\arraybackslash}p{#1}}
\newcolumntype{C}[1]{>{\centering\arraybackslash}p{#1}}
\newcolumntype{R}[1]{>{\raggedleft\arraybackslash}p{#1}}


%more options with enum items
\usepackage{enumitem}

%figures
\usepackage{graphicx}
%image path
\graphicspath{{figures/}}

%acronyms
\usepackage[printonlyused]{acronym} % falls gewünscht kann die Option footnote eingefügt werden, dann wird die Erklärung nicht inline sondern in einer Fußnote dargestellt

% pdf hyperref
\usepackage[
    bookmarks=true,
    bookmarksopen=true,
    bookmarksnumbered=true,
    bookmarksopenlevel=1,
    pdftitle={\titel},
    pdfauthor={\autor},
    pdfcreator={pdflatex},
    pdfsubject={\titel},
    pdfkeywords={\keywords},
    pdflang=German,
    pdfpagelabels=true,
    colorlinks=true,
	linkcolor=LinkColor, 	% Farbe festlegen
	citecolor=LinkColor,
	filecolor=LinkColor,
	menucolor=LinkColor,
	urlcolor=LinkColor,
    plainpages=false,
    hypertexnames=true,
    linktocpage=true,
]{hyperref}

% biblatex setup
% this template uses a custom citation style (wmaainf.bst)
% such a file is not compatible with biblatex and biber 
% thats why biblatex with biber isn't used
% use bibtex instead

%%%%%%%%%%%%%%%%%%%%%%%%%%%%%% configure your listings style %%%%%%%%%%%%%%%%%%%%%%%%%%%%%%%%%%%%%%%%%%%%%%%%
\usepackage{listings}

%%%%%%%%%%%%%%%%%%java settings (default setting)%%%%%%%%%%%%%%%%%%%%%%
\lstset{%
	language=Java,			% default programming language
	numbers=left,			% line numbers left
	stepnumber=1,			% every line is numbered
	numbersep=5pt,			% space beetween line numbers and programm language
	numberstyle=\tiny,		% number style
	breaklines=true,		% line breaking possible
	breakautoindent=true,	% indent after line break
	postbreak=\space,		% space is line break
	tabsize=2,				% tabsize
	basicstyle=\ttfamily\footnotesize, %Monospaced, small font
	showspaces=false,		% show spaces
	showstringspaces=false,	% show spaces in strings 
	extendedchars=true,		% show all characters from the character set Latin1
	captionpos=b,			% sets the caption-position to bottom
	backgroundcolor=\color{ListingBackground}, % background for listing
	xleftmargin=0pt,		 
	xrightmargin=0pt,		
	frame=single,			% frame on
	frameround=ffff,
	rulecolor=\color{darkgray},	% colour of the frame
	fillcolor=\color{ListingBackground},
	keywordstyle=\color{keywords},
	commentstyle=\color{comments},
	stringstyle=\color{strings},
	keywordsprefix = {@}
}
%%%%%%%%%%%%%%%%%%java settings%%%%%%%%%%%%%%%%%%%%%%

%%%%%%%%%%%%%%%%%%yaml as language%%%%%%%%%%%%%%%%%%%%%%
\newcommand\YAMLcolonstyle{\color{red}\mdseries}
\newcommand\YAMLkeystyle{\color{black}\small}
\newcommand\YAMLvaluestyle{\color{black}\footnotesize}

\makeatletter

% here is a macro expanding to the name of the language
% (handy if you decide to change it further down the road)
\newcommand\language@yaml{yaml}

\expandafter\expandafter\expandafter\lstdefinelanguage
\expandafter{\language@yaml}
{
	keywords={true,false,null,y,n},
	keywordstyle=\color{darkgray}\bfseries,
	basicstyle=\YAMLkeystyle, % assuming a key comes first
	sensitive=false,
	comment=[l]{\#},
	morecomment=[s]{/*}{*/},
	commentstyle=\color{purple},
	stringstyle=\YAMLvaluestyle\ttfamily,
	moredelim=[l][\color{orange}]{\&},
	moredelim=[l][\color{magenta}]{*},
	moredelim=**[il][\YAMLcolonstyle{:}\YAMLvaluestyle]{:},   % switch to value style at :
	morestring=[b]',
	morestring=[b]",
	literate =    {---}{{\ProcessThreeDashes}}3
	{>}{{\textcolor{red}\textgreater}}1     
	{|}{{\textcolor{red}\textbar}}1 
	{\ -\ }{{\mdseries\ -\ }}3,
}

% switch to key style at EOL
\lst@AddToHook{EveryLine}{\ifx\lst@language\language@yaml\YAMLkeystyle\fi}
\makeatother

\newcommand\ProcessThreeDashes{\llap{\color{cyan}\mdseries-{-}-}}
%%%%%%%%%%%%%%%%%%yaml as language%%%%%%%%%%%%%%%%%%%%%%

%%%%%%%%%%%%%%%%%%Groovy as language%%%%%%%%%%%%%%%%%%%%

\lstdefinelanguage{Groovy}[]{Java}{
	keywordstyle=\color{keywords},
	stringstyle=\color{strings},
	commentstyle=\color{comments}
	keywords=[3]{each, findAll, groupBy, collect, inject, eachWithIndex},
	morekeywords={def, as, in, use},
	moredelim=[is][\textcolor{darkgray}]{\%\%}{\%\%},
	moredelim=[il][\textcolor{darkgray}]{§§}
}
%%%%%%%%%%%%%%%%%%Groovy as language%%%%%%%%%%%%%%%%%%%%

%%%%%%%%%%%%%%%%%%%%%%%%%%%%%% configure your listings style %%%%%%%%%%%%%%%%%%%%%%%%%%%%%%%%%%%%%%%%%%%%%%%%

% page setup
% \setlength{\topskip}{\ht\strutbox}
\geometry{paper=a4paper,left=2.5cm,top=3.0cm,bindingoffset=.8cm}
\onehalfspacing
\frenchspacing
\clubpenalty = 10000
\widowpenalty = 10000 
\displaywidowpenalty = 10000

% some commands
\newcommand{\ua}{\mbox{u.\,a.\ }}
\newcommand{\zB}{\mbox{z.\,B.\ }}
\newcommand{\dahe}{\mbox{d.\,h.,\ }}
\newcommand{\bzw}{\mbox{bzw.\ }}
\newcommand{\bzgl}{\mbox{bzgl.\ }}
\newcommand{\eg}{\mbox{e.\,g.\ }}
\newcommand{\ie}{\mbox{i.\,e.\ }}
\newcommand{\wrt}{\mbox{w.\,r.\,t.\ }}
\newcommand{\etal}{\mbox{\emph{et.\,al.\ }}}


\usepackage{blindtext}

% load glossary entries
\makenoidxglossaries
\loadglsentries{glossary}

\begin{document}

\setcounter{secnumdepth}{3}  % numerate subsections
\setcounter{tocdepth}{2}  % ...but don't include them in toc

\frontmatter
\thispagestyle{empty}
\pdfbookmark[1]{Cover}{cov}

\begin{titlepage}

\begin{center}

\includegraphics[width=380pt]{figures/ohm-logo.png}\\[1cm]
\LARGE{Fakultät Informatik}\\[2cm]

\huge
\textbf{\titel}\\[1cm]
%
\Large
\artderarbeit~im Studiengang \studiengang\\[1cm]
%
\large
vorgelegt von

\Large
\autor\\[0.5cm]
\small
Matrikelnummer \matrikelnr\\[2cm]

\vspace*{\fill}

\large
\begin{tabular}{p{3cm}p{8cm}}\\
Erstgutachter:  & \quad \erstgutachter\\[1.2ex]
Zweitgutachter: & \quad \zweitgutachter\\[1.2ex]
%discomment "Betreuer" and "Unternehmen" for a thesis in a company
%Betreuer: & \quad \betreuer\\
%Unternehmen: & \quad \unternehmen
\end{tabular}
\end{center}

\begin{center}
\copyright\,\the\year
\end{center}

\vspace{-0.5cm}
\singlespacing
\small
\noindent Dieses Werk einschließlich seiner Teile ist \textbf{urheberrechtlich geschützt}.
Jede Verwertung außerhalb der engen Grenzen des Urheberrechtgesetzes ist ohne Zustimmung des Autors unzulässig und strafbar.
Das gilt insbesondere für Vervielfältigungen, Übersetzungen, Mikroverfilmungen sowie die Einspeicherung und Verarbeitung in elektronischen Systemen.

\end{titlepage}
\cleardoublepage

% TODO %%%%%%%%%%%%%%%%%%%!!!! important for master thesis !!!! %%%%%%%%%%%%%%%%%%%%%%%%%%%%%%%%%%%%%%%%%%%%%%%%%%%%%%%%%%%%%%
% download the following form and complete it (hit save in your editor)
% https://intern.ohmportal.de/fileadmin/Gelenkte_Doks/Abt/SZS/SB/SB_0050_FO_Pruefungsrechtliche_Erklaerung_und_Erklaerung_zur_Veroeffentlichung_der_Abschlussarbeit_public.pdf
%\includepdf{SB_0050_FO_Pruefungsrechtliche_Erklaerung_und_Erklaerung_zur_Veroeffentlichung_der_Abschlussarbeit_public.pdf}\cleardoublepage
%%%%%%%%%%%%%%%%%%%%%%%%%%%%%%%%%%%%%%%%%%%%%%%%%%%%%%%%%%%%%%%%%%%%%%%%%%%%%%%%%%%%%%%%%%%%%%%%%%%%%%%%%%%%%%%%%%%%%%%%

\include{content/0_abstract}\cleardoublepage

\tableofcontents

\mainmatter
\chapter{Introduction}\label{ch:intro}

You may have read about similar things in \cite[S.~X]{Goodliffe2007}.
\par
You can also write footnotes.\footnote{Footnotes will be positioned automatically.}
\blindtext

\blindtext[4]

\section{This is an Important Section}

It is possible to reference glossary entries as \gls{library} as an example.
\blindtext

\subsection{And an even more important subsection}

This shows the use of \ac{API}.
\blindtext

\chapter{Data}\label{ch:data}

\Blindtext

\begin{figure}
	\centering
	\includegraphics[width=200pt]{TH-Nuernberg-RGB.png}\\
	\caption[shortname in list of figures]{caption with label \label{fig: CloudStack}}
\end{figure}

\begin{table}[h]
	\centering
	\begin{tabular}{|C{5cm}|C{8cm}|}
		\hline 
		\textbf{Kopfzeile Spalte 0} & \textbf{Kopfzeile Spalte 1} \\ 
		\hline 
		Zeile 1 Spalte 0 & Zeile 1 Spalte 1 \\
		\hline
		Zeile 2 Spalte 0 & Zeile 2 Spalte 1 \\
		\hline
	\end{tabular}
	\caption[shortname in list of tabels]{\centering caption with labels \label{tab: Erfuellung}}
\end{table}
\chapter{Method}\label{ch:method}

In this chapter, we're actually using some code!

\begin{lstlisting}[language=Python,caption={This is an example of a listing with code in tex},captionpos=b,label=code: code in tex]
x = 1
if x == 1:
    # indented four spaces
    print("x is 1.")

\end{lstlisting}

You can also include listings from a file directly:

\lstinputlisting[language=Python,caption={This is an example of a listing from a file},captionpos=b, label=code: code from file]{listings/example.py}

This was already told in \autoref{ch:intro}.

\begin{itemize}
	\item One Point
	\item The other point
\end{itemize}

\begin{enumerate}
	\item First
	\item Second
\end{enumerate}

This \emph{word} ist emphasized.
\par
Inline Listing is also possible: \lstinline|Set<Words> findDistinct();|
\par
As you can see in \autoref{code: code from file}
\chapter{Experiments}\label{ch:experiments}

\Blindtext

\begin{sidewaysfigure}
	\includegraphics[width=\textwidth]{TH-Nuernberg-RGB}
	\caption[short text]{long caption}
\end{sidewaysfigure}
\include{content/5_outlook}
\include{content/6_summary}

% remove if not needed
\appendix
\chapter{Anhang}\label{app:supplemental-information}

\Blindtext


\chapter{Anderer Anhang}\label{app:supplemental-information}

\Blindtext



\backmatter
\listoffigures
\cleardoublepage

\listoftables
\cleardoublepage

\renewcommand{\lstlistlistingname}{Quelltextverzeichnis}  % change for German thesis
\lstlistoflistings
\cleardoublepage

\bibliographystyle{wmaainf}
\bibliography{refs}

\printnoidxglossaries

%acronyms
%!TEX root = thesis.tex

\chapter{Abkürzungsverzeichnis}
%nur verwendete Akronyme werden letztlich im Abkürzungsverzeichnis des Dokuments angezeigt
%Verwendung: 
%		\ac{Abk.}   --> fügt die Abkürzung ein, beim ersten Aufruf wird zusätzlich automatisch die ausgeschriebene Version davor eingefügt bzw. in einer Fußnote (hierfür muss in header.tex \usepackage[printonlyused,footnote]{acronym} stehen) dargestellt
%		\acs{Abk.}   -->  fügt die Abkürzung ein
%		\acf{Abk.}   --> fügt die Abkürzung UND die Erklärung ein
%		\acl{Abk.}   --> fügt nur die Erklärung ein
%		\acp{Abk.}  --> gibt Plural aus (angefügtes 's'); das zusätzliche 'p' funktioniert auch bei obigen Befehlen
%	siehe auch: http://golatex.de/wiki/%5Cacronym
%	
\begin{acronym}
	
	%A
	\acro{API}{Application-Programming-Interface}
	%B
	%C
	%D
	%E
	%F
	%G
	%H
	%I
	\acro{ISO}{Internationale Organisation für Normung}
	\acroplural{ISO}[ISO]{Internationalen Organisation für Normung}
	%J
	%K
	%L
	%M
	%N
	%O
	%P
	%Q
	%R
	%S
	%T
	%U
	%V
	%W
	%X
	%Y
	%Z
		
\end{acronym}




\end{document}
